\section{Research Questions}

\footnotesize{\textit{Supervisor’s comment (Devon): Research questions should be framed as scientific hypotheses that are falsifiable, precise to your contribution, and informative whether true or false (e.g latent space of medical images are densely packed). Both positive and negative outcomes must yield insights (e.g., if compositional diffusion fails to represent silico-TB, the failure itself teaches us about the latent space structure of medical data). The phrasing should abstract away from narrow cases (e.g., silicosis only) toward principles generalizable across diseases and models. Each hypothesis should be testable with clear validation pathways, and structured such that even its contrapositive or “failure” generates publishable findings.}} \normalsize

\subsection{RQ1: Generative Fidelity of Compositional CXR Synthesis}
\label{subsec:rq1}
\textit{How realistic and clinically valid are chest x-rays generated by compositional diffusion models when simulating silico-tuberculosis compared to single-disease radiographs?}

\footnotesize{\textit{Devon’s comment: Emphasize “clinical validity” as not only realism, but also medical interpretability for radiologists (e.g If we have a nodule from TB and scarring Silicosis, then what does that mean and how is that represented.}} \normalsize

\subsection{RQ2: Faithful Representation of Comorbidity (Silico--TB)}
\label{subsec:rq2}
\begin{enumerate}
    \item \textit{Does the compositional generative process preserve the geometric and pathological integrity of chest x-ray when composing silico-TB?}
    \item \textit{To what extent can compositional diffusion capture the unique interaction of localized TB lesions and diffuse silicosis patterns in a way that conventional augmentation or simple mixing techniques cannot?}
\end{enumerate}

\footnotesize{\textit{Devon’s comment: Stress the hypothesis that synthetic data should reduce brittleness in classifiers when facing atypical or mixed cases.}} \normalsize


\subsection{RQ3: Downstream Utility for Clinical Classifiers}
\label{subsec:rq3}
\begin{enumerate}
    \item \textit{To what extent does incorporating synthetic silico-TB data improve classifier generalization and robustness when faced with atypical or mixed-pathology chest X-rays?}
    \item \textit{Can synthetic silico-TB examples expand the decision space of classifiers so that they are less brittle when confronted with mixed or atypical presentations?}
\end{enumerate}

